%%%%%%%%%%%%%%%%%%%%%%%%%%%%%%%%%%%%
% Lesson Plan (50 minutes)
%%%%%%%%%%%%%%%%%%%%%%%%%%%%%%%%%%%%
\begin{frame}
    \frametitle{Lesson Plan}
    \begin{itemize}
        \item 2 min Lecture: Same Data, Two Lenses --- reintroduce the Lecture 34 drug trial as a case study for both toolkits
        \item 5 min Lecture: The Frequentist Side --- point estimate/CI (and why the CI is invalid here), hypothesis test (and why its conditions fail)
        \item 4 min Lecture: The Bayesian Side --- posterior/credible interval, Bayesian hypothesis test
        \item 4 min Lecture: Side by Side --- comparison table, why the two approaches differed here
        \item 3.5 min Lecture: What If We Had More Data? --- revisit Lecture 31's convergence claim, $n=6$ vs.\ $60$ vs.\ $600$ table
        \item 3 min Lecture: Case Study: Overdispersed Counts --- motivating question, checking the Poisson mean/variance assumption
        \item 6 min Lecture: The Poisson MLE, Fit and Checked; the Negative Binomial distribution; NB shapes for varying $r$ --- so students can see what $r$ does before estimating it themselves
        \item 6 min Think/Pair/Share: \textit{Diagnosing Overdispersion} worksheet --- estimate $r$ for the likes data by method of moments, before the lecture reveals the numerically-optimized $\hat r$ (see \texttt{Lecture35\_worksheet.tex})
        \item 10 min Lecture: rest of Overdispersed Counts case study --- NB MLE via \texttt{optim()}, bootstrap simulation for frequentist uncertainty, MCMC for the Bayesian posterior, side-by-side convergence, what carried over
        \item 2 min Lecture: Module 5 Summary
    \end{itemize}
\end{frame}

%%%%%%%%%%%%%%%%%%%%%%%%%%%%%%%%%%%%
% Total: ~45.5 min content + ~4.5 min buffer for agenda/LO/announcements and pacing slack = 50 min.
% R demo cut (Sec 3.5's convergence table already shows the answer); this is still the densest
% lecture in the module (22 content frames + the TPS activity), so keep an eye on pacing.




